\documentclass[11pt]{article}

\usepackage[margin=0.67in]{geometry}
\usepackage{setspace}
\usepackage{parskip}
\usepackage{hyperref}
\usepackage{titlesec}
\usepackage{xcolor}
\usepackage{array}
\usepackage{ragged2e}

\definecolor{PrimaryRed}{HTML}{8B1E1E}
\definecolor{SoftBlue}{HTML}{DDEEFF}
\definecolor{DarkGray}{HTML}{333333}

\hypersetup{
    colorlinks=true,
    linkcolor=black,
    urlcolor=black,
    citecolor=black
}

\titleformat{\section}{\large\bfseries\color{PrimaryRed}}{}{0em}{}
\titleformat{\subsection}{\normalsize\bfseries}{\thesubsection}{0.5em}{}

\setstretch{1.15}

\begin{document}

\begin{titlepage}
\newgeometry{margin=0.7in}
\thispagestyle{empty}

\begin{center}
    \vspace*{0.25in}
    {\large University of Rochester\\}
    {\normalsize Warner School of Education\\[0.35in]}

    {\color{PrimaryRed}\rule{\textwidth}{1.4pt}}\\[0.22in]
    {\Huge\bfseries Joy-Oriented Literacy Pursuit\\[0.12in]}
    {\Large Formal Reflective Journal\\[0.22in]}
    {\color{PrimaryRed}\rule{0.72\textwidth}{0.7pt}}\\[0.32in]

    \colorbox{SoftBlue}{%
        \parbox{0.86\textwidth}{%
            \centering\vspace{0.12in}
            \textbf{Classroom Technology, Design, and Literacy as Joy}\\[0.05in]
            A Teaching Neurodivergent Journal focused on Literature \& Joy via a Puzzle Plan
            \vspace{0.12in}
        }%
    }
\end{center}

\vspace{0.75in}

\begin{center}
\renewcommand{\arraystretch}{1.35}
\begin{tabular}{>{\bfseries\RaggedLeft}p{1.55in} p{4.85in}}
Student & Piter Garcia \\
Course & EDU498: Culturally-Historically Responsive Literacies \\
Class Session & Module 3 / Week 5 \\
Assignment & Joy-Oriented Literacy Pursuit \\
Context & Teaching Reflection through the Lens of this course \& a Puzzle Plan\\
Submission Date & June 05, 2026 \\
\end{tabular}
\end{center}

\vfill

\begin{center}
\begin{minipage}{0.7\textwidth}
\centering
\itshape
Joy-oriented literacy is not extra. It is part of access. When students feel capable, connected, and allowed to participate as themselves, literacy becomes a tool for liberation, self-determination, and empowerment.
\end{minipage}
\end{center}

\restoregeometry
\end{titlepage}
\clearpage

\doublespacing

\section*{What I Created and Who I Shared It With}

For this joy-oriented literacy pursuit, I designed and facilitated technology lessons where students experienced literacy through building, collaboration, speaking, movement, and problem-solving. The most joyful examples came from Kindergarten circuit and airplane-building lessons and Minecraft lessons in third and fourth grade. What made them literacy pursuits was not the technology itself but the fact that students built, helped, responded, revised, and learned from one another in real time.

In the Kindergarten circuit lessons, students gathered around a shared demonstration, named parts, held materials, and connected pieces before counting down together to turn on a fan. They then worked in teams to build an airplane---assigning roles like wings, body, wheels, and fan connection---so that every student had a meaningful way to participate. In the fourth grade Minecraft lessons, students who finished early helped classmates who were still working. That shift turned task completion into community responsibility and made peer teaching part of the lesson itself.

\section*{My Creative and Learning Process}

My process began with a question: \textit{how can a technology lesson become more than task completion?} Knobel and Lankshear (2014) helped me answer it: literacy is not fixed to a single form; it evolves as tools, communities, and purposes change. Price-Dennis, Holmes, and Smith (2015) pushed further by showing that digital tools only become inclusive when teachers intentionally use them to reduce barriers, increase participation, and support meaningful communication. Together, their work reframed my lessons as genuine literacy work---students were interpreting systems, building vocabulary, communicating with peers, and representing thinking through more than one mode.

In the circuit and airplane activities, students connected vocabulary to action. They did not only hear the word ``switch''; they held it, connected it, and watched what happened when it worked. Role distribution made the work collaborative rather than individual. In Minecraft, inviting early finishers to help peers shifted the room from ``I am done'' to ``I can help someone else succeed''---a moment where one student troubleshoots while another consolidates mastery.

Joy was not an accident of fun activities. It happened because students had access to materials, roles, language, peer support, and visible success. Knobel and Lankshear describe new literacies as social practices built around participation and community. Price-Dennis and colleagues argue that adding a device to a traditional task is not innovation; the real question is whether the task creates more access, voice, and critical thinking. That became the design principle: \textit{enough structure to feel safe, enough freedom to feel ownership}.

\section*{How This Cultivated Joy, Love, and Aesthetic Fulfillment}

Joy came from students seeing their work become real. In the Kindergarten fan-circuit lessons, the class counted down together and watched the fan turn on. The science was no longer abstract---it was visible, audible, and shared. Aesthetic fulfillment came from building something with shape, motion, and purpose: their work had form, effort, and a relationship to the physical world.

Love came through relationships. In the fourth grade Minecraft lessons, students who understood the task helped classmates who were still working. Success was no longer private or competitive; it became something students could give each other.

Joy in literacy does not always look quiet or neat. Sometimes it looks like a student trying again, a group arguing over what goes where, a fan finally turning on, or a peer stepping in to help. That is aesthetic fulfillment: learning feels alive, students can see themselves inside the work, and the classroom becomes a place where knowledge is built rather than delivered.

\section*{Muhammad's Framework and the Case for Literacy as Access}

Muhammad's framework asks that a literacy pursuit develop identity, skills, intellect, criticality, and joy together. This project did not aim only at science skills; it asked students to see themselves as builders, helpers, and problem-solvers---because literacy should help students understand themselves, participate in community, and experience agency.

Knobel and Lankshear (2014) gave me language for what I was already seeing: building, explaining, and peer teaching are new literacies because meaning-making happens through digital, visual, participatory, and networked forms, not only through print. Price-Dennis, Holmes, and Smith (2015) connect this to inclusion: when tasks allow for multimodal creation, oral explanation, and collaborative design, more students can show what they know.

Brooks and Frankel (2020) matter here because students do not arrive as blank technology users. Their work on transnational digital literacy shows that students may already be navigating languages, countries, and communities through digital tools before school names any of it academic. Building in Minecraft, explaining a circuit orally, helping a peer troubleshoot---these are literacy acts that require interpretation, cultural knowledge, and audience awareness.

Melo-Pfeifer and Gertz (2022) pushed me toward criticality: digital spaces are not neutral, and students need tools to question headlines, missing voices, representation, and assumptions. The Common Sense Digital Citizenship materials extend this by giving students language for privacy, confirmation bias, online identity, and upstander communication. My Puzzle Plan builds exactly those habits: collaboration, questioning, audience awareness, and responsibility to others.

Stornaiuolo and Thomas (2017) reframe students as active participants who use digital tools to speak and create public meaning---not passive consumers. Pinkard's ecosystem view reinforces this: students need tools, mentors, audiences, and real reasons to create. The Grade 4 peer-teaching moment was that ecosystem in action. Serpagli and Mensah (2021) connect to my placement directly: students explained technical ideas through familiar, social, and active forms. Across all these readings, the argument is consistent---digital literacy positions students as creators.

\section*{How This Informs My Future Teaching}

Literacy is not limited to written language. It is also building, coding, designing, explaining, translating, troubleshooting, and helping. My classroom will offer multiple entry points so students can show understanding through writing, speaking, drawing, recording, building, or peer teaching---especially for students who are multilingual, neurodivergent, disabled, culturally diverse, or different from the ``default'' learner schools often imagine.

I will keep designing for authentic audiences: Minecraft worlds tied to a science concept, bilingual multimodal explanations, collaborative murals, peer-created tutorials, or circuitry projects documented through video or drawing. Peer teaching will be built into the culture, not treated as a bonus, because students often understand each other's barriers better than adults do.

The core question my Puzzle Plan and EQUITAS lens keeps asking is not whether a task looks traditional enough to count as literacy. It is whether the task gives students real access to meaning, identity, criticality, and agency. Knobel and Lankshear show that literacy changes with communities and tools; Price-Dennis and colleagues show that inclusion is designed, not assumed; Brooks and Frankel show that students' existing practices already are literacy; and Stornaiuolo and Thomas show that digital tools can be instruments of voice. When students feel capable, connected, and allowed to participate as themselves, literacy becomes a tool for liberation, self-determination, and empowerment.

\end{document}
